\section{Experimental Setup and Benchmarking}
\label{sec:experimental_setup}

Having established the theoretical design and optimisation of the \texttt{DetectorNX} architecture, both the \gls{snn} and its \gls{cnn} counterpart,
the final step is to formulate a rigorous evaluation strategy to empirically validate its performance against the continuous \gls{cnn} baseline.
In accordance with the project's primary objectives, the evaluation cannot be limited to a singular metric of bounding box accuracy; rather, it must holistically assess the neuromorphic paradigm.
To achieve this, a comprehensive benchmarking suite was designed (as mentioned in Figure \ref{fig:system_architecture_block}) covering seven distinct experimental protocols:
\begin{itemize}
    \item Dynamic Power Estimation
    \item Quantisation Error Analysis (Accuracy vs. Object Size)
    \item Confidence Calibration
    \item Temporal Evidence Accumulation
    \item Real-Time Inference Latency
    \item Theoretical Energy Consumption
    \item Global Evaluator: Mean \gls{iou}, \gls{map_50}, Recall, and Average Detections Per Frame
\end{itemize}

\subsection{Dynamic Power Estimation}
\label{subsec:dynamic_power}

A fundamental objective of this research is to quantify the energy efficiency gains achieved by transitioning from a continuous \gls{cnn} to a bio-plausible \gls{snn}.
Traditional \gls{cnn} architectures are characterised by static power consumption; they perform a dense, fixed number of \gls{mac} operations for every forward pass, regardless of the information density within the input scene. 
In contrast, \glspl{snn} leverage event-driven dynamics, where metabolic energy is consumed only when a neuron reaches its firing threshold and generates a discrete spike.

The Dynamic Power experiment investigates the reactive scalability of the architecture by mapping instantaneous energy consumption against the temporal density of the input stream. 
This analysis serves to validate the event-driven nature of the \gls{snn}, demonstrating that its metabolic cost is a dynamic function of activation sparsity—mimicking biological neural efficiency by consuming energy only in response to informative stimuli.

\subsubsection{\gls{cnn} Power Estimation}

In a standard \gls{cnn}, the total energy per frame ($E_{CNN}$) is calculated by multiplying the total number of \glspl{flop} by the cost of a 32-bit \gls{mac} ($E_{MAC}$):

\begin{equation}
  E_{CNN} = \sum_{l=1}^{L} (FLOPs_l \times E_{MAC})
\end{equation}

\noindent Where $E_{MAC} = 4.6\text{ pJ}$ per 32-bit floating point operation, as referenced in Table \ref{tab:energy_constants} \cite{horowitz20141}.

The number of \glspl{flop} is the total number of \glspl{mac} operations required to produce the output feature map; in a \gls{cnn}, this number is static, regardless of what the image contains.
The \glspl{flop} of a \gls{cnn} can be calculated as follows:

\begin{equation}
     FLOPs_l = (H_{out} \times W_{out}) \times (K \times K) \times C_{in} \times C_{out}
\end{equation}

Where $(H_{out} \times W_{out})$ is the number of pixels in the output frame, $(K \times K)$ is the size of the convolutional kernel, $C_{in}$ is the number of input channels, and $C_{out}$ is the number of output channels.


\subsubsection{\gls{snn} Power Estimation}

For the spiking architecture, the calculation is scaled by the activation sparsity ($\zeta$) of each layer.
The metabolic cost of the spiking backbone can be quantified via the \gls{sop} metric;
unlike the static complexity of the \gls{cnn} baseline, the \gls{sop} count for a given layer is a dynamic function, which can be calculated as follows:
\begin{equation}
SOPs_l = FLOPs_l \times \zeta \times T
\label{eq:sop_zeta}
\end{equation}

\noindent where $\zeta_l$ represents the average firing rate of layer $l$, $T$ is the number of temporal simulation steps ($T=10$).

In order to obtain $\zeta_l$, a specialised \texttt{SpikeRateMonitor} tool was developed.
The tool attaches a hook to every \gls{lif} inside the \gls{snn} backbone, which listens to the layer during the forward pass and intercepts the binary spike tensor before it moves to the next layer.
For every layer $l$, the \gls{snn} produces a spike tensor $S_l$ with the shape $[T, B, C, H, W]$, where $T = 10$ is the number of timesteps, $B$ is the batch size, and $C, H, W$ are the spatial dimensions.
Since a spike is binary ($1$ for fire, $0$ for rest), $\zeta$ can be calculated by taking the global mean of the entire tensor:

\begin{equation}
    \zeta_l = \frac{1}{N} \sum_{i=1}^{N} S_{l, i}
\end{equation}

To calculate Dynamic Power, $\zeta$ was calculated by running inference over all validation samples and taking the mean firing rate per layer across all samples in the dataset.

% \todo[inline]{Mention SpikeRateMonitor in Evaluation Suite Section and Diagram}

Because \glspl{snn} utilise binary activations, the expensive multiplication part of the \gls{mac} operation is omitted, leaving only a \gls{ac} of the weight to the membrane potential.
The total energy per frame ($E_{SNN}$) is defined as:

\begin{equation}
\label{eq:e_snn}
  E_{SNN} = \sum_{l=1}^{L} (SOPs_l \times E_{AC})
\end{equation}


\noindent where $E_{AC} = 0.9\text{ pJ}$ is the energy cost of a single floating-point addition, as referenced in Table \ref{tab:energy_constants} \cite{horowitz20141}.


\subsection{Quantization Error Analysis (Accuracy vs. Object Size)}
\label{subsec:quantization_error}

The Quantization Error Analysis experiment investigates the discretization threshold of the \gls{snn} architecture and quantifies the extent to which the 1-bit binary nature of spiking activations limits the model's geometric acuity. 
Unlike the continuous \gls{cnn} baseline, which utilizes 32-bit floating-point activations to represent spatial coordinates, the \gls{snn} is constrained by the finite temporal resolution of its simulation window ($T=10$).

\subsubsection{Core Metrics and Scale Partitioning}
The primary performance metric for this experiment is the \gls{miou}, which quantifies the overlap between the predicted bounding box ($P$) and the ground truth ($GT$). To analyze the impact of scale on regression precision, the validation set is partitioned into three distinct categories based on the absolute pixel area ($A_{px}$) of the ground truth objects:
\begin{itemize}
    \item \textbf{Small:} Objects in the bottom 33rd percentile of area (typically distant vehicles or small artifacts).
    \item \textbf{Medium:} Objects in the middle 33rd percentile.
    \item \textbf{Large:} Objects in the top 33rd percentile (vehicles in close proximity to the sensor).
\end{itemize}

\subsubsection{Mathematical Framework}
The categorization and evaluation are governed by two fundamental geometric formulas. First, the pixel area is calculated by denormalizing the ground truth dimensions relative to the raw sensor resolution:
\begin{equation}
    A_{px} = (W_{norm} \cdot W_{sensor}) \times (H_{norm} \cdot H_{sensor})
\end{equation}
where $W_{sensor} = 1280$ and $H_{sensor} = 720$. Subsequently, the precision of the model's spatial regression is quantified via the \gls{iou} formula:
\begin{equation}
    IoU = \frac{\text{Area}(P \cap GT)}{\text{Area}(P \cup GT)}
\end{equation}

\subsubsection{Hypothesis: The Scale-Dependent Accuracy Gap}
We hypothesize a significant scale-dependent accuracy gap between the two architectures.
For large-scale objects, the \gls{snn} is expected to achieve parity with the \gls{cnn} baseline: because the coordinate values are high, the minor fluctuations inherent in binary spike-trains represent a negligible percentage of the total object area, allowing for stable regression.

Conversely, for small-scale objects, a substantial performance divergence is anticipated. 
In an \gls{snn} with $T=10$, a neuron can only represent $T+1$ discrete states ($0/10, 1/10, \dots, 10/10$). 
As documented by \cite{cordone2022objectdetectionspikingneural}, this rounding error, or quantization penalty, is massive relative to the dimensions of small vehicles. 
If a target requires a precise coordinate gradient (e.g., $0.154$) for accurate localization, the \gls{snn} is physically forced to discretize this value, establishing a fundamental ``resolution floor'' that limits \gls{miou} on small features. 
While the use of \texttt{SEWResBlock} architectures improves signal flow \cite{fang2022resnet_sew}, it does not inherently resolve this quantization of the final regression output, which is further exacerbated by the low event counts typical of small objects in event-based vision \cite{gallego2020event}.

\subsection{Confidence Calibration}
\label{subsec:confidence_calibration}

The Confidence Calibration experiment evaluates predictive reliability by investigating the calibration gap between spiking and continuous architectures. Modern \glspl{cnn} are prone to systematic overconfidence, frequently assigning high objectness scores to inaccurate detections \cite{guo2017calibration}. 
Conversely, the temporal dynamics of \glspl{snn} may offer superior uncertainty quantification, providing a more robust reliability profile for safety-critical automotive tasks \cite{skatchkovsky2022bayesian}.

\subsubsection{Methodology: Reliability Diagrams}
To quantify calibration, we generate a Reliability Diagram mapping predicted confidence ($C$) against actual regression accuracy ($IoU$) for every validation detection.
In a perfectly calibrated system, these values align along the identity line ($C = IoU$). 
Clumping in the high-confidence, low-accuracy quadrant indicates an overconfidence bias, where the model fails to signal its own predictive errors.

\subsubsection{Hypothesis: Spikes as Uncertainty Sensors}
The \gls{cnn} baseline is hypothesised to exhibit significant miscalibration due to its continuous activation functions and exact gradient descent, which tend to overfit the confidence head. 
The \gls{snn} is expected to demonstrate a more linear, calibrated trend. 
In the spiking domain, low-quality features generate fewer spikes, leading to lower membrane potentials in the detection head. 
This mechanism allows the \gls{snn} to function as an implicit uncertainty sensor, where discrete spike rates effectively gate the output confidence logit relative to the strength of the temporal evidence.

\subsection{Temporal Evidence Accumulation}
\label{subsec:temporal_evidence}

The Temporal Evidence Accumulation experiment investigates the sequential decision-making dynamics inherent to the \gls{snn} architecture. 
Unlike the \gls{cnn} baseline, which performs a single spatial pass over the fully integrated $100\text{ms}$ voxel, the \gls{snn} processes the event stream iteratively across $T=10$ timesteps. 
This property makes \glspl{snn} intrinsically suited to low-latency spatiotemporal integration --- the capacity 
to produce meaningful predictions before the full observation window has elapsed \cite{pfeiffer_snn_opportunities}. 
The experiment therefore examines whether this theoretical advantage manifests empirically: specifically, whether the 
model can produce a valid bounding box regression before the $100\text{ms}$ window is complete.

In the context of vehicle object detection, this is a highly significant experiment --- for example,
if the \gls{snn} can accurately show a bounding box at $t=5$ ($50\text{ms}$) instead of the full $t=10$ ($100\text{ms}$) as the \gls{cnn} would,
this would lead to a 50\% reduction in detection time. In a time-sensitive domain such as vehicle and pedestrian detection, such advantages could be the difference between a crash and a safe stop.

\subsubsection{Methodology: Temporal Snapshots}
To evaluate the evolution of predictions, we extract intermediate outputs at three specific checkpoints: $t=2$ (20ms), $t=5$ (50ms), and $t=10$ (100ms). 
For each snapshot, the candidate bounding box coordinates and instantaneous confidence scores are recorded to observe the \textit{hunting behavior} of the regression head:
the observable process where the networks' predicted bounding box moves, vibrates or searches for the object's true edges before finally settling into a stable, high-precision detection. 
This process is driven by the discrete \gls{lif} dynamics:

\begin{equation}
    v_i[t] = v_i[t-1] + \frac{1}{\tau} \left( \sum_j w_{ij}s_j[t] - (v_i[t-1] - V_{reset}) \right)
\end{equation}

where $v_i[t]$ represents the accumulated evidence (membrane potential) and $\sum w_{ij}s_j[t]$ represents the incoming observations (spikes). 
This recursive filtering allows the network to update its belief about vehicle localization with every incoming pulse of data.

\subsubsection{Hypothesis: Bounding Box Sharpening}
% We hypothesize a phenomenon of \textit{Bounding Box Sharpening} over time. At $t=20\text{ms}$, we anticipate low confidence and high coordinate jitter, as most \gls{lif} neurons will not have aggregated sufficient potential to cross the firing threshold ($V_{th}=0.2$). As $t$ progresses toward 100ms, the integration of further spike-trains should allow the model to filter out stochastic noise and consolidate the structured geometric edges of the vehicle. This transition from early-stage estimation to high-precision localization mirrors the ``rapid categorization'' observed in the human visual system, where initial spikes are sufficient to drive meaningful, albeit lower-precision, semantic decisions \cite{thorpe1996speed}.

We hypothesize a phenomenon of \textit{Bounding Box Sharpening} over time. 
At $t = 20\text{ms}$, we anticipate high coordinate variance, as the majority 
of \gls{lif} neurons will not have accumulated sufficient membrane potential to 
cross the firing threshold ($V_{th} = 0.2$) within the first two of ten 
timesteps. As $t$ progresses toward $100\text{ms}$, continued spike train 
integration allows stochastic membrane noise to average out, consolidating 
responses to the vehicle's geometric edges into a stable coordinate estimate.

This progressive refinement is analogous in spirit to the ``rapid categorisation'' 
phenomenon documented in the human visual system, wherein early feedforward spikes 
drive coarse semantic decisions that sharpen as further neural evidence accumulates 
\cite{thorpe1996speed}. While Thorpe et al.\ address categorical recognition rather 
than coordinate regression, the underlying principle --- that sparse early spike 
activity produces meaningful but imprecise representations which sharpen over time 
--- is consistent with the anticipated behaviour of the \gls{snn} backbone.
\subsection{Real-Time Latency Benchmark}
\label{subsec:exp_latency_setup}

This experiment quantifies the forward-pass inference speed of both architectures
on conventional \gls{gpu} hardware at batch size $B=1$, simulating a live 
single-stream event camera deployment. Latency ($\text{ms/f}$) and throughput 
($\text{FPS} = 1000 / \text{ms/f}$) are measured using \texttt{torch.cuda.Event} 
timing to capture true \gls{gpu} execution time, bypassing \gls{cpu}-to-\gls{gpu} 
communication overhead.

\subsubsection{Hypothesis}
The \gls{cnn} baseline is expected to exhibit lower latency, as \glspl{gpu} are 
optimised for the dense matrix multiplications that underpin its single flattened 
spatial pass. The \gls{snn} must instead simulate \gls{lif} membrane dynamics 
iteratively:

\begin{equation}
    v[t] = v[t-1] + x[t] - \beta \cdot v[t-1]
    \label{eq:lif_update_latency}
\end{equation}

\noindent This constitutes a hardware-mismatch penalty: the latency advantage of 
\glspl{snn} only fully materialises on asynchronous neuromorphic hardware, not on 
synchronous \gls{gpu} accelerators \cite{davies2018loihi, blouw2019benchmarking}. 
The experiment aims to quantify this penalty and assess whether \texttt{step\_mode=`m'} 
parallelisation \cite{fang2023spikingjelly} brings the \gls{snn} within a 
$2\times$--$3\times$ latency margin of the \gls{cnn}.

\subsubsection{Limitations}
The measured \gls{gpu} latency figures cannot be extrapolated to neuromorphic 
silicon. The A100's $312\ \text{TFLOPS}$ (FP16) characterises dense matrix 
throughput, which cannot be judged against asynchronous spike-based execution. 
While Loihi benchmarks report up to $10\times$ lower latency than embedded
\glspl{gpu} for batch-size-1 inference \cite{davies2018loihi}, these figures are
task-specific; direct deployment on neuromorphic hardware is identified as future
work.
\subsection{Theoretical Energy Consumption}
\label{subsec:energy_consumption}


\subsection{Global Performance Evaluator}
\label{subsec:global_evaluator}

The Global Performance Evaluator serves as the primary comparative evaluation between the \gls{snn} and \gls{cnn} architectures,
aggregating performance across the full 200-sample validation set.
Unlike the preceding targeted experiments, this benchmark establishes the \textit{Performance-Efficiency Frontier}:
quantifying the accuracy penalty incurred by transitioning to discrete spiking dynamics, whilst validating whether the multi-box refinements and others in the \texttt{ULTRA} revision of \texttt{DetectorNX}
mentioned in Section \ref{sec:architectural_evolution} are sufficient to resolve scene-level detection ambiguity.

\subsubsection{Metrics}

In this experiment, four complementary metrics are evaluated:

\begin{itemize}
    \item \textbf{\gls{miou}:} measures geometric precision between predicted and ground truth bounding boxes.

    % \item \textbf{Recall:} measures the proportion of ground truth objects successfully localised
    % at confidence threshold $\tau = 0.5$, penalising architectures that fail to respond to
    % sparse event regions %\cite{gallego2020event}.%

    \item \textbf{Recall:}
    measures the proportion of ground truth objects successfully localised.

\end{itemize}

% \textbf{Recall} measures the proportion of ground truth objects successfully localised
% at confidence threshold $\tau = 0.5$, penalising architectures that fail to respond to
% sparse event regions \parencite{gallego2020event}. \textbf{Average Detections per Frame
% (ADP)} quantifies detector selectivity relative to the ground truth mean of 2.02 objects
% per frame; significant positive deviation indicates spurious ``ghost'' detections arising
% from background noise misclassification. \textbf{Mean Prediction Confidence} captures
% network self-certainty by averaging sigmoid-activated confidence scores across all
% positive detections, probing whether spiking dynamics produce decisive or hesitant class
% activations \parencite{eshraghian2021training}.

% \subsubsection*{Hypothesis}

% We hypothesise that whilst the CNN establishes the upper mIoU bound, the SNN will
% achieve \textit{detection parity} in recall and ADP, yielding an average detection
% count consistent with the ground truth distribution. This would demonstrate that
% spiking temporal dynamics are sufficient for objectness discrimination without
% continuous-valued feature maps \parencite{roy2019towards,
% cordone2022objectdetectionspikingneural}. A measurable confidence gap is also
% anticipated: the SNN's binary spike outputs are expected to produce lower but temporally
% stable confidence distributions across the $T = 10$ timestep window, consistent with
% prior observations that surrogate-trained networks distribute representational energy
% across timesteps rather than single forward passes \parencite{fang2022resnet_sew}.

% \subsubsection*{Significance}

% This experiment provides three contributions. It delivers the \textit{parity proof},
% framing the accuracy-efficiency tradeoff quantitatively. It offers \textit{logical
% validation} that the Phase 1 ambiguity problem has been resolved through the
% architectural refinements described in Section~\ref{sec:architectural_evolution}.
% Finally, evaluation across the full validation set ensures statistical robustness,
% confirming that results are representative of the broader ETraM environment rather
% than artefacts of favourable individual frames.